%% -*- mode: LaTeX; compile-command: "lhs2TeX --poly explaining-errors.lhs -o explaining-errors.tex && pdflatex explaining-errors.tex" -*-
\documentclass[xcolor=svgnames,12pt,aspectratio=169]{beamer}

%% ODER: format ==         = "\mathrel{==}"
%% ODER: format /=         = "\neq "
%
%
\makeatletter
\@ifundefined{lhs2tex.lhs2tex.sty.read}%
  {\@namedef{lhs2tex.lhs2tex.sty.read}{}%
   \newcommand\SkipToFmtEnd{}%
   \newcommand\EndFmtInput{}%
   \long\def\SkipToFmtEnd#1\EndFmtInput{}%
  }\SkipToFmtEnd

\newcommand\ReadOnlyOnce[1]{\@ifundefined{#1}{\@namedef{#1}{}}\SkipToFmtEnd}
\usepackage{amstext}
\usepackage{amssymb}
\usepackage{stmaryrd}
\DeclareFontFamily{OT1}{cmtex}{}
\DeclareFontShape{OT1}{cmtex}{m}{n}
  {<5><6><7><8>cmtex8
   <9>cmtex9
   <10><10.95><12><14.4><17.28><20.74><24.88>cmtex10}{}
\DeclareFontShape{OT1}{cmtex}{m}{it}
  {<-> ssub * cmtt/m/it}{}
\newcommand{\texfamily}{\fontfamily{cmtex}\selectfont}
\DeclareFontShape{OT1}{cmtt}{bx}{n}
  {<5><6><7><8>cmtt8
   <9>cmbtt9
   <10><10.95><12><14.4><17.28><20.74><24.88>cmbtt10}{}
\DeclareFontShape{OT1}{cmtex}{bx}{n}
  {<-> ssub * cmtt/bx/n}{}
\newcommand{\tex}[1]{\text{\texfamily#1}}	% NEU

\newcommand{\Sp}{\hskip.33334em\relax}


\newcommand{\Conid}[1]{\mathit{#1}}
\newcommand{\Varid}[1]{\mathit{#1}}
\newcommand{\anonymous}{\kern0.06em \vbox{\hrule\@width.5em}}
\newcommand{\plus}{\mathbin{+\!\!\!+}}
\newcommand{\bind}{\mathbin{>\!\!\!>\mkern-6.7mu=}}
\newcommand{\rbind}{\mathbin{=\mkern-6.7mu<\!\!\!<}}% suggested by Neil Mitchell
\newcommand{\sequ}{\mathbin{>\!\!\!>}}
\renewcommand{\leq}{\leqslant}
\renewcommand{\geq}{\geqslant}
\usepackage{polytable}

%mathindent has to be defined
\@ifundefined{mathindent}%
  {\newdimen\mathindent\mathindent\leftmargini}%
  {}%

\def\resethooks{%
  \global\let\SaveRestoreHook\empty
  \global\let\ColumnHook\empty}
\newcommand*{\savecolumns}[1][default]%
  {\g@addto@macro\SaveRestoreHook{\savecolumns[#1]}}
\newcommand*{\restorecolumns}[1][default]%
  {\g@addto@macro\SaveRestoreHook{\restorecolumns[#1]}}
\newcommand*{\aligncolumn}[2]%
  {\g@addto@macro\ColumnHook{\column{#1}{#2}}}

\resethooks

\newcommand{\onelinecommentchars}{\quad-{}- }
\newcommand{\commentbeginchars}{\enskip\{-}
\newcommand{\commentendchars}{-\}\enskip}

\newcommand{\visiblecomments}{%
  \let\onelinecomment=\onelinecommentchars
  \let\commentbegin=\commentbeginchars
  \let\commentend=\commentendchars}

\newcommand{\invisiblecomments}{%
  \let\onelinecomment=\empty
  \let\commentbegin=\empty
  \let\commentend=\empty}

\visiblecomments

\newlength{\blanklineskip}
\setlength{\blanklineskip}{0.66084ex}

\newcommand{\hsindent}[1]{\quad}% default is fixed indentation
\let\hspre\empty
\let\hspost\empty
\newcommand{\NB}{\textbf{NB}}
\newcommand{\Todo}[1]{$\langle$\textbf{To do:}~#1$\rangle$}

\EndFmtInput
\makeatother
%
%
%
%
%
%
% This package provides two environments suitable to take the place
% of hscode, called "plainhscode" and "arrayhscode". 
%
% The plain environment surrounds each code block by vertical space,
% and it uses \abovedisplayskip and \belowdisplayskip to get spacing
% similar to formulas. Note that if these dimensions are changed,
% the spacing around displayed math formulas changes as well.
% All code is indented using \leftskip.
%
% Changed 19.08.2004 to reflect changes in colorcode. Should work with
% CodeGroup.sty.
%
\ReadOnlyOnce{polycode.fmt}%
\makeatletter

\newcommand{\hsnewpar}[1]%
  {{\parskip=0pt\parindent=0pt\par\vskip #1\noindent}}

% can be used, for instance, to redefine the code size, by setting the
% command to \small or something alike
\newcommand{\hscodestyle}{}

% The command \sethscode can be used to switch the code formatting
% behaviour by mapping the hscode environment in the subst directive
% to a new LaTeX environment.

\newcommand{\sethscode}[1]%
  {\expandafter\let\expandafter\hscode\csname #1\endcsname
   \expandafter\let\expandafter\endhscode\csname end#1\endcsname}

% "compatibility" mode restores the non-polycode.fmt layout.

\newenvironment{compathscode}%
  {\par\noindent
   \advance\leftskip\mathindent
   \hscodestyle
   \let\\=\@normalcr
   \let\hspre\(\let\hspost\)%
   \pboxed}%
  {\endpboxed\)%
   \par\noindent
   \ignorespacesafterend}

\newcommand{\compaths}{\sethscode{compathscode}}

% "plain" mode is the proposed default.
% It should now work with \centering.
% This required some changes. The old version
% is still available for reference as oldplainhscode.

\newenvironment{plainhscode}%
  {\hsnewpar\abovedisplayskip
   \advance\leftskip\mathindent
   \hscodestyle
   \let\hspre\(\let\hspost\)%
   \pboxed}%
  {\endpboxed%
   \hsnewpar\belowdisplayskip
   \ignorespacesafterend}

\newenvironment{oldplainhscode}%
  {\hsnewpar\abovedisplayskip
   \advance\leftskip\mathindent
   \hscodestyle
   \let\\=\@normalcr
   \(\pboxed}%
  {\endpboxed\)%
   \hsnewpar\belowdisplayskip
   \ignorespacesafterend}

% Here, we make plainhscode the default environment.

\newcommand{\plainhs}{\sethscode{plainhscode}}
\newcommand{\oldplainhs}{\sethscode{oldplainhscode}}
\plainhs

% The arrayhscode is like plain, but makes use of polytable's
% parray environment which disallows page breaks in code blocks.

\newenvironment{arrayhscode}%
  {\hsnewpar\abovedisplayskip
   \advance\leftskip\mathindent
   \hscodestyle
   \let\\=\@normalcr
   \(\parray}%
  {\endparray\)%
   \hsnewpar\belowdisplayskip
   \ignorespacesafterend}

\newcommand{\arrayhs}{\sethscode{arrayhscode}}

% The mathhscode environment also makes use of polytable's parray 
% environment. It is supposed to be used only inside math mode 
% (I used it to typeset the type rules in my thesis).

\newenvironment{mathhscode}%
  {\parray}{\endparray}

\newcommand{\mathhs}{\sethscode{mathhscode}}

% texths is similar to mathhs, but works in text mode.

\newenvironment{texthscode}%
  {\(\parray}{\endparray\)}

\newcommand{\texths}{\sethscode{texthscode}}

% The framed environment places code in a framed box.

\def\codeframewidth{\arrayrulewidth}
\RequirePackage{calc}

\newenvironment{framedhscode}%
  {\parskip=\abovedisplayskip\par\noindent
   \hscodestyle
   \arrayrulewidth=\codeframewidth
   \tabular{@{}|p{\linewidth-2\arraycolsep-2\arrayrulewidth-2pt}|@{}}%
   \hline\framedhslinecorrect\\{-1.5ex}%
   \let\endoflinesave=\\
   \let\\=\@normalcr
   \(\pboxed}%
  {\endpboxed\)%
   \framedhslinecorrect\endoflinesave{.5ex}\hline
   \endtabular
   \parskip=\belowdisplayskip\par\noindent
   \ignorespacesafterend}

\newcommand{\framedhslinecorrect}[2]%
  {#1[#2]}

\newcommand{\framedhs}{\sethscode{framedhscode}}

% The inlinehscode environment is an experimental environment
% that can be used to typeset displayed code inline.

\newenvironment{inlinehscode}%
  {\(\def\column##1##2{}%
   \let\>\undefined\let\<\undefined\let\\\undefined
   \newcommand\>[1][]{}\newcommand\<[1][]{}\newcommand\\[1][]{}%
   \def\fromto##1##2##3{##3}%
   \def\nextline{}}{\) }%

\newcommand{\inlinehs}{\sethscode{inlinehscode}}

% The joincode environment is a separate environment that
% can be used to surround and thereby connect multiple code
% blocks.

\newenvironment{joincode}%
  {\let\orighscode=\hscode
   \let\origendhscode=\endhscode
   \def\endhscode{\def\hscode{\endgroup\def\@currenvir{hscode}\\}\begingroup}
   %\let\SaveRestoreHook=\empty
   %\let\ColumnHook=\empty
   %\let\resethooks=\empty
   \orighscode\def\hscode{\endgroup\def\@currenvir{hscode}}}%
  {\origendhscode
   \global\let\hscode=\orighscode
   \global\let\endhscode=\origendhscode}%

\makeatother
\EndFmtInput
%


\usepackage[all]{xy}
\usepackage{brent}
\usepackage{xspace}
\usepackage{fancyvrb}
\usepackage{ulem}

\usepackage{ucs}
\usepackage[utf8x]{inputenc}
% \usepackage[backend=pgf,extension=pgf,input,outputdir=diagrams]{diagrams-latex}
\graphicspath{{images/}}

\usepackage{wasysym}  %% for \frownie

%%%%%%%%%%%%%%%%%%%%%%%%%%%%%%%%%%%%%%%%%%%%%%%%%%%%%%%%%%%%
%%%%%%%%%%%%%%%%%%%%%%%%%%%%%%%%%%%%%%%%%%%%%%%%%%%%%%%%%%%%

% Math typesetting

%% a bit more space for matrices
\setlength{\arraycolsep}{5pt}

\newcommand{\ty}[3]{{#1} \vdash {#2} : {#3}}
\newcommand{\nty}[3]{{#1} \nvdash {#2} : {#3}}

%%%%%%%%%%%%%%%%%%%%%%%%%%%%%%%%%%%%%%%%%%%%%%%%%%%%%%%%%%%%
%%%%%%%%%%%%%%%%%%%%%%%%%%%%%%%%%%%%%%%%%%%%%%%%%%%%%%%%%%%%

\newcommand{\etc}{\textit{etc.}}
\renewcommand{\eg}{\textit{e.g.}\xspace}
\renewcommand{\ie}{\textit{i.e.}\xspace}

\newcommand{\theschool}{Off the Beaten Track}
% \newcommand{\thelocation}{Los Angeles}
\newcommand{\thedate}{13 January 2018}

%%%%%%%%%%%%%%%%%%%%%%%%%%%%%%%%%%%%%%%%%%%%%%%%%%%%%%%%%%%%
%%%%%%%%%%%%%%%%%%%%%%%%%%%%%%%%%%%%%%%%%%%%%%%%%%%%%%%%%%%%

\setbeamertemplate{items}[circle]

\mode<presentation>
{
  \usetheme{default}                          % use a default (plain) theme

  \setbeamertemplate{navigation symbols}{}    % don't show navigation
                                              % buttons along the
                                              % bottom
  \setbeamerfont{normal text}{family=\sffamily}

  % XX remove this before giving actual talk!
  % \setbeamertemplate{footline}[frame number]
  % {%
  %   \begin{beamercolorbox}{section in head/foot}
  %     \vskip2pt
  %     \hfill \insertframenumber
  %     \vskip2pt
  %   \end{beamercolorbox}
  % }

  \AtBeginSection[]
  {
    \begin{frame}<beamer>
      \frametitle{}

      \begin{center}
%        \includegraphics[width=2in]{\sectionimg}
%        \bigskip

        {\Huge \insertsectionhead}
      \end{center}
    \end{frame}
  }
}

\defbeamertemplate*{title page}{customized}[1][]
{
  \vbox{}
  \vfill
  \begin{centering}
    \begin{beamercolorbox}[sep=8pt,center,#1]{title}
      \usebeamerfont{title}\inserttitle\par%
      \ifx\insertsubtitle\@empty%
      \else%
        \vskip0.25em%
        {\usebeamerfont{subtitle}\usebeamercolor[fg]{subtitle}\insertsubtitle\par}%
      \fi%
    \end{beamercolorbox}%
    \vskip1em\par
    {\usebeamercolor[fg]{titlegraphic}\inserttitlegraphic\par}
    \vskip1em\par
    \begin{beamercolorbox}[sep=8pt,center,#1]{author}
      \usebeamerfont{author}\insertauthor
    \end{beamercolorbox}
    \begin{beamercolorbox}[sep=8pt,center,#1]{institute}
      \usebeamerfont{institute}\insertinstitute
    \end{beamercolorbox}
    \begin{beamercolorbox}[sep=8pt,center,#1]{date}
      \usebeamerfont{date}\insertdate
    \end{beamercolorbox}
  \end{centering}
  \vfill
}

\newenvironment{xframe}[1][]
  {\begin{frame}[fragile,environment=xframe,#1]}
  {\end{frame}}

% uncomment me to get 4 slides per page for printing
% \usepackage{pgfpages}
% \pgfpagesuselayout{4 on 1}[uspaper, border shrink=5mm]

% \setbeameroption{show only notes}

\renewcommand{\emph}{\textbf}

\title{Explaining Type Errors}
\date{\theschool \\ \thedate}
\author{\usebeamercolor[fg]{title}{Brent Yorgey} \and Richard
  Eisenberg \and Harley Eades}
% \titlegraphic{\includegraphics[width=1in]{deriv-tree}}

%%%%%%%%%%%%%%%%%%%%%%%%%%%%%%%%%%%%%%%%%%%%%%%%%%%%%%%%%%%%

\begin{document}

\begin{xframe}{}
   \titlepage
\end{xframe}

\begin{xframe}{The Dreaded Type Error Message}
\footnotesize
\begin{Verbatim}
Could not deduce (Num t0)
from the context: (Num (t -> a), Num t, Num a)
  bound by the inferred type for 'it':
             forall a t. (Num (t -> a), Num t, Num a) => a
  at <interactive>:4:1-19
The type variable 't0' is ambiguous
In the ambiguity check for the inferred type for 'it'
To defer the ambiguity check to use sites, enable AllowAmbiguousTypes
When checking the inferred type
  it :: forall a t. (Num (t -> a), Num t, Num a) => a
\end{Verbatim}
\end{xframe}

\begin{xframe}{}
  \begin{center}
    \includegraphics{what-the-function.jpg}
  \end{center}
\end{xframe}

\begin{xframe}{Theses}
  \begin{itemize}
  \item ``Improving'' error messages doesn't fundamentally help.
  \item Interactive \emph{error explanations} instead of static
    \emph{error messages}.
  \item \emph{Error explanation} = \emph{constructive evidence} for
    an error.
  \end{itemize}
\end{xframe}

\section{The Curse of Information}

\begin{xframe}{}
  \begin{center}
\begin{BVerbatim}
(\f -> f 3) (\p -> fst p)
\end{BVerbatim}
  \end{center}
\end{xframe}

\begin{xframe}{}
%  $\app{(\uabs{f}{\app{f}{3}})}{(\uabs{p}{\app{fst}{p}})}$

  \begin{center}
    \begin{BVerbatim}
(\f -> f 3) (\p -> fst p)
    \end{BVerbatim}

    \begin{Verbatim}
Type mismatch between expected type (t, b0) and actual type Int
    \end{Verbatim}
  \end{center}
\end{xframe}

\begin{xframe}
%  $\app{(\uabs{f}{\app{f}{3}})}{(\uabs{p}{\app{fst}{p}})}$

  \begin{center}
    \begin{BVerbatim}
(\f -> f 3) (\p -> fst p)
    \end{BVerbatim}

\begin{Verbatim}
Type mismatch between expected type (t, b0) and actual type Int
  In the first argument of fst, namely p
  In the expression: fst p
  In the first argument of \ f -> f 3, namely
    (\ p -> fst p)
\end{Verbatim}
  \end{center}
\end{xframe}

\begin{xframe}
  \footnotesize
\begin{Verbatim}
    Type mismatch between expected type (t, b0) and actual type Int
      In the first argument of fst, namely p
      In the expression: fst p
      In the first argument of \ f -> f 3, namely
        (\ p -> fst p)
      Relevant bindings include:
        fst :: (a,b) -> a
      Inferred types for subterms:
        3             :: Int
        (\f -> f 3)   :: forall a. (Int -> a) -> a
        (\p -> fst p) :: (Int -> a0)
        (\f -> f 3) (\p -> fst p) :: a0
      Suggested fixes:
        Change p to (p,y)
        Change fst to a function expecting an Int
        Change 3 to (x,y)
\end{Verbatim}
\end{xframe}

% XXX IF TIME: make one more slide, even more ridiculously small &
% crowded with info

\begin{xframe}{}
  \begin{center}
    \includegraphics[width=3in]{not-helping.jpg}
  \end{center}
\end{xframe}

\begin{xframe}
  \footnotesize
  \begin{Verbatim}
    Type mismatch between expected type (t, b0) and actual type Int
      In the first argument of fst, namely p
      In the expression: fst p
      In the first argument of \ f -> f 3, namely
        (\ p -> fst p)
      Relevant bindings include:
        fst :: (a,b) -> a
      Inferred types for subterms:
        3             :: Int
        (\f -> f 3)   :: forall a. (Int -> a) -> a
        (\p -> fst p) :: (Int -> a0)
        (\f -> f 3) (\p -> fst p) :: a0
      Suggested fixes:
        Change p to (p,y)
        Change fst to a function expecting an Int
        Change 3 to (x,y)
  \end{Verbatim}
\end{xframe}

\begin{xframe}{}
  \Large
  \begin{center}
    \sout{MESSAGES} \\
    $\Downarrow$ \\
    EXPLANATIONS
  \end{center}
\end{xframe}

\begin{xframe}{}
  \small
\begin{Verbatim}
p is expected to have a pair type but was inferred to have type Int.
  + Why is p expected to have a pair type?
  + Why was p inferred to have type Int?







\end{Verbatim}
\end{xframe}

\begin{xframe}{}
  \small
\begin{Verbatim}
p is expected to have a pair type but was inferred to have type Int.
  + Why is p expected to have a pair type?
  - Why was p inferred to have type Int?
    => p is the parameter of the lambda expression \p -> fst p, which
       must have type (Int -> a0).
    + Why must (\p -> fst p) have type (Int -> a0)?




\end{Verbatim}
\end{xframe}

\begin{xframe}{}
  \small
\begin{Verbatim}
p is expected to have a pair type but was inferred to have type Int.
  + Why is p expected to have a pair type?
  - Why was p inferred to have type Int?
    => p is the parameter of the lambda expression \p -> fst p, which
       must have type (Int -> a0).
    - Why must (\p -> fst p) have type (Int -> a0)?
      => It is an argument to (\f -> f 3), which was inferred to have
         type (forall a. (Int -> a) -> a).
      + Why was (\f -> f 3) inferred to have type
        (forall a. (Int -> a) -> a)?
\end{Verbatim}
\end{xframe}

\begin{xframe}{Related work\dots}
  \begin{itemize}
  \item Plociniczak \& Odersky: Scalad (2012)
  \item Stuckey, Sulzmann \& Wazny: Chameleon (2003)
  \item Simon, Chitil, \& Huch (2000)
  \item Beaven \& Stansifer (1993)
  \end{itemize}

  Also\dots
  \begin{itemize}
  \item Seidel, Jhala \& Weimer: dynamic witnesses for errors
    (ICFP 2016). A great idea, but not the focus of this talk.
  \end{itemize}
\end{xframe}
\section{Explaining errors}

\begin{xframe}{The type of type inference?}

  \begin{overprint}
\onslide<1>
\begin{hscode}\SaveRestoreHook
\column{B}{@{}>{\hspre}l<{\hspost}@{}}%
\column{E}{@{}>{\hspre}l<{\hspost}@{}}%
\>[B]{}\Varid{infer}\mathbin{:}\Conid{Context}\to \Conid{Term}\to \Conid{Maybe}\;\Conid{Type}{}\<[E]%
\ColumnHook
\end{hscode}\resethooks

\onslide<2>
\begin{hscode}\SaveRestoreHook
\column{B}{@{}>{\hspre}l<{\hspost}@{}}%
\column{E}{@{}>{\hspre}l<{\hspost}@{}}%
\>[B]{}\Varid{infer}\mathbin{:}\Conid{Context}\to \Conid{Term}\to \Conid{Maybe}\;\Conid{TypingDerivation}{}\<[E]%
\ColumnHook
\end{hscode}\resethooks

\onslide<3>
\begin{hscode}\SaveRestoreHook
\column{B}{@{}>{\hspre}l<{\hspost}@{}}%
\column{E}{@{}>{\hspre}l<{\hspost}@{}}%
\>[B]{}\Varid{infer}\mathbin{:}\Conid{Context}\to \Conid{Term}\to (\Conid{Error}\mathbin{+}\Conid{TypingDerivation}){}\<[E]%
\ColumnHook
\end{hscode}\resethooks

\onslide<4>
\begin{hscode}\SaveRestoreHook
\column{B}{@{}>{\hspre}l<{\hspost}@{}}%
\column{E}{@{}>{\hspre}l<{\hspost}@{}}%
\>[B]{}\Varid{infer}\mathbin{:}\Conid{Context}\to \Conid{Term}\to (\neg\;\Conid{TypingDerivation}\mathbin{+}\Conid{TypingDerivation}){}\<[E]%
\ColumnHook
\end{hscode}\resethooks

\onslide<5->
\begin{hscode}\SaveRestoreHook
\column{B}{@{}>{\hspre}l<{\hspost}@{}}%
\column{E}{@{}>{\hspre}l<{\hspost}@{}}%
\>[B]{}\Varid{infer}\mathbin{:}\Conid{Context}\to \Conid{Term}\to (\Conid{UntypingDerivation}\mathbin{+}\Conid{TypingDerivation}){}\<[E]%
\ColumnHook
\end{hscode}\resethooks

\end{overprint}

\begin{overprint}
  \onslide<6>
  See Ulf Norell invited talk @ ICFP 2013: \\
  \url{http://www.cse.chalmers.se/~ulfn/code/icfp2013/ICFP.html}

  \onslide<7->
  \begin{center}
    To generate interactive error explanations, \\
    \emph{focus on designing untyping derivations}.
  \end{center}
\end{overprint}

\onslide<8>
\begin{center}
  (or unparsing derivations, or nontermination derivations, or \dots?)
\end{center}
\end{xframe}

\begin{xframe}{Example: STLC + \N}
\begin{align*}
  t &::= x \mid n \mid t_1 + t_2 \mid \abs x \tau t \mid \app{t_1}{t_2} \\
  \tau &::= \mathbb{N} \mid \tau_1 \to \tau_2 \\
  \Gamma &::= \varnothing \mid \Gamma,x:\tau
\end{align*}
\end{xframe}

\begin{xframe}{Example: STLC + \N}
\framebox{$\ty \Gamma t {\tau}$}
\begin{mathpar}
  \inferrule{x : \tau \in \Gamma}{\ty \Gamma x \tau} \and
  \inferrule{\ty {\Gamma, x:\tau_1} t {\tau_2}}{\ty \Gamma {\abs x
      {\tau_1} t}{\tau_1 \to \tau_2}} \and
  \inferrule{\ty \Gamma {t_1} {\tau_1 \to \tau_2} \\ \ty \Gamma {t_2} {\tau_1}}
      {\ty \Gamma {\app{t_1}{t_2}} {\tau_2}} \\
  \inferrule{ }{\ty \Gamma n \N} \and
  \inferrule{\ty \Gamma {t_1}{\N} \\ \ty \Gamma {t_2}{\N}}{\ty \Gamma
    {t_1 + t_2}{\N}}
\end{mathpar}
\end{xframe}

\begin{xframe}{Untyping for STLC + $\N$}
  \framebox{$\nty \Gamma t {\tau}$}
  \begin{overprint}
    \onslide<1>
    \begin{mathpar}
      \inferrule*[right=Mismatch\frownie{}]{\ty \Gamma t {\tau_1} \\ \tau_1 \neq \tau_2}{\nty
        \Gamma t {\tau_2}}
    \end{mathpar}

    \onslide<2>
    \begin{mathpar}
      \inferrule{ }{\N \neq (\tau_1 \to \tau_2)} \and
      \inferrule{\tau_1 \neq \tau_2}{(\tau_1 \to \tau_3) \neq (\tau_2
        \to \tau_4)} \and \dots
    \end{mathpar}

    \onslide<3>
    \begin{mathpar}
      \inferrule*[right=PlusL\frownie{}]{\nty \Gamma {t_1} {\N}}{\nty \Gamma {t_1 + t_2}
        {\tau}} \and
      \inferrule*[right=PlusR\frownie{}]{\nty \Gamma {t_2} {\N}}{\nty \Gamma {t_1 + t_2}
        {\tau}} \and
      \inferrule*[right=PlusTy\frownie{}]{\tau \neq \N}{\nty \Gamma {t_1 + t_2} {\tau}}
    \end{mathpar}

    \onslide<4>
    \begin{mathpar}
      \inferrule*[right=AbsTy\frownie{}]
      {\forall \tau_2.\; \tau \neq (\tau_1 \to \tau_2)}
      {\nty \Gamma {\abs x {\tau_1} t} \tau}
      \and

  % -- Otherwise, τ is of the form (τ₁ ⇒ τ₂) but the body t does not
  % -- have type τ₂.  Note this could be either because t is not typable
  % -- at all, or because it has some type other than τ₂.
  % ƛ        : ∀ {n} {Γ : Ctx n} {t} {τ₁ τ₂}
  %            → (τ₁ ∷ Γ) ⊬ t ∶ τ₂
  %            → Γ ⊬ ƛ τ₁ t ∶ (τ₁ ⇒ τ₂)
      \inferrule*[right=AbsBody\frownie{}]
      {\nty{\Gamma, x:\tau_1} {t} {\tau_2}}
      {\nty \Gamma {\abs x {\tau_1} t} {\tau_1 \to \tau_2}}
    \end{mathpar}

    \onslide<5>
    \begin{mathpar}
      % ·-fun    : ∀ {n} {Γ : Ctx n} {t₁ t₂} {τ₂}
      % → (∀ {τ₁} → Γ ⊬ t₁ ∶ τ₁ ⇒ τ₂)
      % → Γ ⊬ t₁ · t₂ ∶ τ₂
      \inferrule*[right=LhsTy\frownie{}]
      {\forall \tau_1.\; \nty \Gamma {t_1} {\tau_1 \to \tau_2}}
      {\nty \Gamma {\app {t_1} {t_2}} {\tau_2}}
      \and

      % ·-arg    : ∀ {n} {Γ : Ctx n} {t₁ t₂} {τ₁ τ₂}
      % → Γ ⊢ t₁ ∶ τ₁ ⇒ τ₂
      % → Γ ⊬ t₂ ∶ τ₁
      % → Γ ⊬ t₁ · t₂ ∶ τ₂
      \inferrule*[right=RhsTy\frownie{}]
      {\ty \Gamma {t_1} {\tau_1 \to \tau_2} \\
        \nty \Gamma {t_2} {\tau_1}}
      {\nty \Gamma {\app {t_1} {t_2}} {\tau_2}}
    \end{mathpar}
  \end{overprint}
\end{xframe}

\begin{xframe}{Example}
  \begin{mathpar}
    \inferrule*[right=AbsBody\frownie{}]
    {
      \inferrule*[right=PlusTy\frownie{}]
      {\N \to \N \neq \N
      }
      {\nty {f : \N \to \N}{f + 2}{\N \to \N}}
    }
    {\nty \varnothing {\abs{f}{\N \to \N}{f + 2}}{(\N \to \N) \to \N
        \to \N}}
  \end{mathpar}
\end{xframe}

\begin{xframe}{Correctness?}
  \[ \neg \ty \Gamma t \tau \iff \nty \Gamma t \tau \] \medskip

  \onslide<2->
  ($\Rightarrow$): follows from decidability of type
  inference/checking. \medskip

  \onslide<3->
  ($\Leftarrow$): more interesting; $\nty \Gamma t \tau$ adequately
  encodes evidence that $t$ does not have type $\tau$. \medskip

  \onslide<4>
  Still a lot of room for variations: round-tripping need not
  be the identity!
\end{xframe}

\begin{xframe}{Exploring explanations}
  \begin{itemize}
  \item<+-> Untyping derivation = typing derivation with one negated
    subterm?
  \item<+-> Untyping derivation = typing derivation with many negated
    subterms?
  \item<+-> Untyping derivation = typing derivation with negated subterms
    pushed as far to the leaves as possible?
  \item<+-> \dots pulled as far toward the root as possible?
  \item<+-> Treat typing rules as a big sum-of-products, negate, and
    apply De Morgan laws?
  \end{itemize}
\end{xframe}

% \begin{xframe}{Unification?}
  
% \end{xframe}

\end{document}
